% ===== latex-recomposed:document =====
% flm: style file flmrecomposed.91ff4b0d.sty
\begin{equation}
  a = b\flmLabels{\label{ref1}}
\end{equation}

See \eqref{ref1}.
%%
% ===== latex-recomposed:style flmrecomposed.91ff4b0d.sty =====
\NeedsTeXFormat{LaTeX2e}
\ProvidesPackage{flmrecomposed.91ff4b0d}
% --- packages ---
% flm.latex.pkg.amsmath
\RequirePackage{amsmath}
% --- commands and environments the library defines for the document ---
% flm.latex.def.labels
% \flmLabels{<\label commands>}: the labels of one place in the document.
% In text it runs them as written. On a line of a display-math
% environment, where amsmath keeps the last label of a line and loses the
% ones before it, it runs the first and records every further one as an alias:
% a line per alias is written to the .aux file after the last page and read on
% the next run, making the alias refer to what the first label refers to
% (for \ref and for \zcref alike). Redefine \flmLabels to change what a
% group of labels does; \def\flmLabels#1{#1} restores plain LaTeX behavior,
% amsmath's "Multiple \label's" message included.
\ProvideDocumentCommand\flmLabels{m}{%
  \ifmmode\expandafter\flm@labels@math\else\expandafter\@firstofone\fi{#1}}
% Reading the group: the first \label as written, the rest as aliases of it.
\def\flm@labels@math#1{\def\flm@lbl@first{}\flm@lbl@scan#1\flm@lbl@stop}
\def\flm@lbl@scan{\@ifnextchar\label\flm@lbl@take\flm@lbl@rest}
\def\flm@lbl@take\label#1{%
  \ifx\flm@lbl@first\@empty
    \def\flm@lbl@first{#1}%
    \label{#1}%
  \else
    \expandafter\flm@recordalias\expandafter{\flm@lbl@first}{#1}%
  \fi
  \flm@lbl@scan}
\def\flm@lbl@rest#1\flm@lbl@stop{#1}
% Collecting the aliases. A name is recorded once however often the line is
% typeset (amsmath measures a line before setting it).
\def\flm@aliaslist{}
\def\flm@recordalias#1#2{%
  \ifcsname flm@alias@#2\endcsname\else
    \expandafter\gdef\csname flm@alias@#2\endcsname{}%
    \g@addto@macro\flm@aliaslist{\flm@alias@entry{#2}{#1}}%
  \fi}
% Writing them, from the hook that runs after the last page has been shipped
% out: every \newlabel of that page has reached the .aux file by then, so the
% alias lines come after them. \AtEndDocument would run too early.
\def\flm@alias@write#1#2{\immediate\write\@auxout{\string\flm@aliaslabel{#1}{#2}}}
\AddToHook{enddocument/afterlastpage}{\let\flm@alias@entry\flm@alias@write\flm@aliaslist}
% Reading them back: the alias refers to what the original refers to. The
% assignments are global because the .aux file is read inside a group.
\def\flm@aliaslabel#1#2{%
  \ifcsname r@#2\endcsname
    \global\expandafter\let\csname r@#1\expandafter\endcsname\csname r@#2\endcsname
  \fi
  \ifcsname Z@R@#2\endcsname
    \global\expandafter\let\csname Z@R@#1\expandafter\endcsname\csname Z@R@#2\endcsname
  \fi}
\endinput
